\subsubsection{Hard Attention to the Task (HAT)}

\gls{hat} \cite{ll_hat_2018} is a parameter-isolation \gls{cl} method that
learns a near-binary attention mask for each task via differentiable task
embeddings. Unlike pruning-based approaches, the masks are trained end-to-end
and harden progressively through a sigmoid annealing schedule, so that capacity
allocation emerges from the optimisation rather than a post-hoc pruning step.

\paragraph{Mask Embeddings and Gate Computation}

For each gated layer $l$ the model maintains an embedding matrix
$\mathbf{E}_l \in \mathbb{R}^{T \times d_l}$, where $T$ is the number of tasks
and $d_l$ is the number of units (channels) in that layer. During training on
task $t$, the gate vector for layer $l$ is
%
\begin{equation}\label{eqn:hat_gate}
    \mathbf{m}_l^{(t)}
    =
    \sigma\!\bigl(s\;\mathbf{e}_l^{(t)}\bigr)
    \in (0,1)^{d_l} ,
\end{equation}
%
where $\mathbf{e}_l^{(t)} = \mathbf{E}_l[t,\,:]$ is the $t$-th row of the
embedding table and $s > 0$ is the annealing scale. Because $d_l$ counts
channels rather than individual weights, the mask operates at
\emph{channel granularity}: one scalar gate per channel, broadcast across all
spatial positions. This is coarser than \gls{packnet}, which selects individual
weights by magnitude, but the channel mask is learned end-to-end rather than
assigned post-hoc. The gate modulates the
layer's output activations element-wise,
%
\begin{equation}\label{eqn:hat_apply}
    \tilde{\mathbf{h}}_l = \mathbf{h}_l \odot \mathbf{m}_l^{(t)} .
\end{equation}

\paragraph{Effective Weight Mask}

Because each weight $W_{ij}^{(l)}$ connects unit $i$ of layer $l{-}1$ to unit
$j$ of layer $l$, its effective mask is determined jointly by the post-gate of
the preceding layer (pre-gate) and the post-gate of the current layer:
%
\begin{equation}\label{eqn:hat_weight_mask}
    u_{ij}^{(l)}
    =
    \min\!\bigl(m_{l-1,\,i}^{(t)},\; m_{l,\,j}^{(t)}\bigr) .
\end{equation}
%
This \emph{min} rule ensures that a weight is considered allocated to task $t$
only if both its input unit and output unit are active under the task mask. In the
codebase this is computed per-tensor by broadcasting the pre- and post-gate
vectors to match the parameter shape.

\paragraph{Mask Annealing}

The scaling parameter $s$ in \eqref{eqn:hat_gate} controls mask sharpness.
Within each task's training period, $s$ is annealed from $1/s_\text{max}$ to
$s_\text{max}$ over $B$ batches. Two schedules are available via the
\texttt{anneal\_schedule} hyperparameter.

The \emph{geometric} schedule follows the original paper \cite{ll_hat_2018},
interpolating log-linearly:
%
\begin{equation}\label{eqn:hat_anneal_geometric}
    s(b)
    =
    s_\text{max}^{\,2b/(B-1)\,-\,1} ,
\end{equation}
%
which passes through $s = 1$ at the midpoint and keeps masks near-linear for
the first half of training before sharpening rapidly toward the end.

The \emph{linear} schedule (codebase default) interpolates arithmetically:
%
\begin{equation}\label{eqn:hat_anneal_linear}
    s(b)
    =
    \frac{1}{s_\text{max}}
    +
    \frac{b}{B-1}
    \left(s_\text{max} - \frac{1}{s_\text{max}}\right) ,
\end{equation}
%
where $b$ is the current batch index. Both schedules start at $1/s_\text{max}$
and end at $s_\text{max}$, but the linear version reaches $s_\text{max}/2$ at
the midpoint and therefore hardens masks substantially earlier.

To compensate for the amplifying effect of $s$ on embedding gradients, the
gradient of each embedding is rescaled after the backward pass,
%
\begin{equation}\label{eqn:hat_grad_comp}
    \nabla_{\mathbf{e}}
    \leftarrow
    \frac{s_\text{max}}{s}
    \cdot
    \frac{\cosh(s\,\mathbf{e}) + 1}
         {\cosh(\mathbf{e}) + 1}
    \odot
    \nabla_{\mathbf{e}} .
\end{equation}
%
The two factors address distinct sources of scale variation. The gradient
flowing back through the gate is proportional to $s\,\sigma'(s\mathbf{e})$.
Using the identity $\sigma'(x) = \bigl[2(\cosh(x)+1)\bigr]^{-1}$, this equals
$s\,/\,\bigl[2(\cosh(s\mathbf{e})+1)\bigr]$. The $s_\text{max}/s$ factor
cancels the explicit $s$. The $\cosh$ ratio then converts
$\sigma'(s\mathbf{e}) \to \sigma'(\mathbf{e})$ by replacing
$(\cosh(s\mathbf{e})+1)$ with $(\cosh(\mathbf{e})+1)$: when $s$ is large the
sigmoid saturates and $\sigma'(s\mathbf{e})$ near-vanishes, while the
denominator $(\cosh(\mathbf{e})+1)$ is independent of $s$. After both
corrections the compensated gradient equals
$s_\text{max}\,\nabla_{\mathbf{m}}\,\sigma'(\mathbf{e})$---equivalent to
fixing $s = 1$ and scaling by $s_\text{max}$---so the embeddings receive a
consistent learning signal throughout annealing. The clamp prevents
$\cosh(s\mathbf{e})$ from overflowing when $s$ is large.

\paragraph{Cumulative Mask and Gradient Protection}

At the end of task $t$ the masks are binarised at $s = s_\text{max}$ and
accumulated into a cumulative importance mask,
%
\begin{equation}\label{eqn:hat_cumask}
    \mathbf{M}_{<t+1}^{(l)}
    =
    \max\!\Bigl(
        \mathbf{M}_{<t}^{(l)},\;
        \mathbf{1}\!\bigl[\mathbf{m}_l^{(t)} \geq \tfrac{1}{2}\bigr]
    \Bigr)
    \in \{0,1\}^{d_l} .
\end{equation}
%
The indicator maps the soft gate $\mathbf{m}_l^{(t)} \in (0,1)^{d_l}$ to a
binary vector, so $\mathbf{M}_{<t+1}^{(l)}$ and all subsequent cumulative masks
are binary. When training on any subsequent task $t' > t$, the corresponding
binary effective weight mask is
%
\begin{equation}\label{eqn:hat_cumweight_mask}
    u_{ij}^{(l,<t')}
    =
    \min\!\bigl(M_{<t',\,i}^{(l-1)},\; M_{<t',\,j}^{(l)}\bigr)
    \in \{0,1\} ,
\end{equation}
%
and gradients are hard-zeroed wherever past allocation is indicated,
%
\begin{equation}\label{eqn:hat_grad_protect}
    \nabla_{W_{ij}^{(l)}}
    \leftarrow
    \nabla_{W_{ij}^{(l)}}
    \cdot
    \bigl(1 - u_{ij}^{(l,<t')}\bigr) .
\end{equation}
%
Because $u_{ij}^{(l,<t')} \in \{0,1\}$, the factor $(1 - u_{ij}^{(l,<t')})$
is also binary: gradients are either fully zeroed or fully passed through.
This prevents gradient updates for future tasks from modifying parameters
important to any past task.

\paragraph{Loss and Capacity Regularisation}

The training objective for task $t$ is
%
\begin{equation}\label{eqn:hat_loss}
    \mathcal{L}_\text{HAT}
    =
    \mathcal{L}_\text{CE}
    +
    \gamma\,R\!\left(\mathbf{M}^{(t)},\,\mathbf{M}_{<t}\right) ,
\end{equation}
%
where $\gamma$ controls regularisation strength. The capacity penalty $R$
encourages the current task's masks to occupy only previously unused units,
%
\begin{equation}\label{eqn:hat_reg}
    R
    =
    \frac{
        \displaystyle\sum_{l,\,j}
        m_{l,j}^{(t)}\!\bigl(1 - M_{<t,j}^{(l)}\bigr)
    }{
        \displaystyle\sum_{l,\,j}
        \bigl(1 - M_{<t,j}^{(l)}\bigr)
    } .
\end{equation}
%
For the first task ($t = 0$), where there is no cumulative mask, $R$ reduces to
the mean of all gate values, penalising unnecessary capacity use from the outset.
Embedding weights are also clamped to $[-6, 6]$ after each step to prevent the
embedding values from saturating the sigmoid.

\paragraph{Inference}

At test time the masks are evaluated at $s = s_\text{max}$, producing near-binary
gates that effectively select a task-specific sub-network.
